% !TeX root = ../neuroforge_cfd.tex
% ---------------------------------------------------------------------------
% Residual-floor section. Self-contained \input section.
% Needs amsmath/amsthm + \newtheorem{theorem} and \newtheorem{proposition}
% (declared in the preamble).
% Cross-refs OUT: sec:iters, tab:iters, sec:descent, sec:selective,
%   sec:floor_ladder, sec:indist-ablation, sec:method.
% Labels IN (referenced from body.tex -- do not rename):
%   sec:residual_floor, thm:residual-floor.
% New labels: thm:consistency-floor, prop:kernel, eq:rf-decomp, eq:rf-lower,
%   eq:omitted-term.
% Citations used: krishnapriyan2021failure, wang2022ntk, trefethenbau1997,
%   beckerrannacher2001dwr, stetter1978defect, brandt2011multigrid,
%   morin2000data, bochevgunzburger2009book, rhiechow1983, lei2026newtonkrylov,
%   zhang2026phymgn.
%
% REWRITE 2026-09-07 (theory/formal-guarantees). Changes from the previous version:
%  - opening sentence corrected: tab:iters shows residual RISING while error FALLS.
%  - NEW Theorem 1 (consistency floor): proves (H2) has a nonzero continuum limit,
%    which body.tex's contribution list already asserts in bold and which existed
%    nowhere in the paper. Includes the exact 2-D identity |T| = sqrt2 ||S||_F |grad nu_t|.
%  - the closure-omission mechanism is DEMOTED throughout: it is ~4.5% of the floor
%    and full stress repair moves the floor <0.1% (floor_resolution_decomposition.json).
%  - detector paragraph no longer rests on sigma_min (the constant purged from leg
%    (iii)); restated as a measured between-case claim.
%  - kernel paragraph replaced by Proposition 1 (counting bound, exact solid
%    invisibility, exact diagonal nu_t block, Nyquist/Rhie-Chow mode).
%  - contribution (a) no longer advertises the upper bound ||r*||/sigma_min.
%  - goal-oriented/operator-availability line added (the novelty sentence).
% ---------------------------------------------------------------------------

\section{The residual floor: what an operator-inconsistent monitor can and cannot do}
\label{sec:residual_floor}

The iteration sweep of \autoref{sec:iters} (\autoref{tab:iters}) shows a dissociation:
under iteration the monitored physics residual \emph{rises} ($0.113\to0.618$) while the
field error \emph{falls} ($3.92\to2.29$, before rising again to $2.58$ at the largest
caps). That is an anti-correlation along one trajectory of one checkpoint, not a causal
claim; the causal experiment---descending $J=\tfrac12\|R_h\|^2$ directly, with no network
in the loop---is \autoref{sec:descent}. This section supplies the mechanism common to
both, and it is not ill-conditioning. It is that the monitored residual is
\emph{self-consistent at a field that is not the truth}: the operator a deployed surrogate
can afford to evaluate is not the operator whose root its training labels are, so the
truth carries a floor and the objective's minimiser sits away from $u^\star$.

The section is organised around the paper's three verbs. \textbf{Certify: only loosely},
and this is the leg the theory settles---\autoref{thm:consistency-floor} shows the floor
survives the continuum limit, and \autoref{thm:residual-floor} shows the truth is neither a
minimiser nor a stationary point, so no \emph{threshold} on the residual norm implies a
bound on the error. A calibrated bound is still available by conformalising the score
against held-out data, and we report one; what the floor destroys is its tightness, since
$86\%$ of a typical prediction's score is present at zero error and the resulting
certificate is $5.6$--$6.8\times$ the quantity being certified (\autoref{sec:conformal}).
\textbf{Descend: no}, by \autoref{thm:residual-floor}(iii)--(iv), measured in
\autoref{sec:descent}. \textbf{Rank: yes}---nothing here contradicts that, and we are
explicit below about why a floor that defeats certification leaves between-case triage
intact, because the same monitor flags worst-decile drag error at AUROC $0.952$
(\autoref{sec:selective}).

\subsection{Setup: the monitored operator, stated exactly}
\label{sec:rf_operator}

Let $\Omega\subset\mathbb{R}^2$ be the crop domain itself (a $3c\times3c$ square about a
unit chord here), $\Omega_h$ the uniform Cartesian grid on it, $h$ its spacing,
$M\subset\Omega_h$ the fluid cells, $W\subset M$ the \emph{wall ring} (fluid cells with a
$4$-neighbour in the solid), and $A:=M\setminus W$ the \textbf{active set}. We write
$\Omega_f\subset\Omega$ for the open fluid region, so $|\Omega_f|<|\Omega|$. The state is
$u=(u,v,p,\nu_t)$. The monitored operator $R_h$ evaluates, on physical (denormalised)
fields with $\nu_{\mathrm{eff}}=\nu+\nu_t$ pointwise,
\begin{equation}
\begin{aligned}
r_c&=D_xu+D_yv,\\
r_x&=uD_xu+vD_yu+D_xp-\nu_{\mathrm{eff}}\Delta_hu,\\
r_y&=uD_xv+vD_yv+D_yp-\nu_{\mathrm{eff}}\Delta_hv,
\end{aligned}
\label{eq:monitored}
\end{equation}
with $D_x,D_y$ second-order central in the interior and first-order one-sided at the two
edges of each axis, and $\Delta_h$ the compact three-point second difference per axis. All
three are set to zero on the solid and on $W$, then non-dimensionalised by $U/\ell$ (continuity)
and $U^2/\ell$ (momentum), with $U$ the freestream speed and $\ell$ the reference length
(we reserve $L$ for the Jacobian below). The scalar we write $\|R_h(\cdot)\|$ throughout is the deployed one,
\texttt{Diagnostics.residual\_norm}: the root-mean-square over the \textbf{whole grid} of
$\sqrt{r_c^2+r_x^2+r_y^2}$ after that non-dimensionalisation, with \textbf{no boundary term}.
We name it because the paper also reports band-restricted norms over a region fixed in
physical units when it studies refinement (\autoref{sec:floor_ladder}); the two are
different scoring regions and their numbers are not interchangeable.

Write $J(u)=\tfrac12\|R_h(u)\|^2$, let $u^\star$ be the discrete ground-truth field, define
the \emph{residual floor} $r^\star:=R_h(u^\star)$, and for a prediction $\hat u$ write
$e:=\hat u-u^\star$ with
\begin{equation}
R_h(\hat u)=r^\star+L\,e+o(\|e\|),\qquad L:=DR_h(u^\star),
\label{eq:rf-decomp}
\end{equation}
where---here and throughout---$L$ denotes the Jacobian restricted to the fluid,
$L:\mathbb{R}^{4|M|}\to\mathbb{R}^{3|A|}$.

\subsection{The floor does not vanish in the continuum limit}

The objection this theorem answers is the one a numerical analyst raises first: \emph{a
discrete operator applied to a non-solution leaves a truncation error; refine the grid.}
It does not apply, because $R_h$ is not a low-order discretisation of the RANS equations.
It is a \emph{consistent} discretisation of a \emph{different} system. Recall that $R_h$ is
consistent with a PDE system $F(u)=0$ if $R_h(u|_{\Omega_h})\to0$ pointwise as $h\to0$ for
every sufficiently smooth solution $u$ of $F$.

\begin{theorem}[Consistency floor: (H2) is a theorem, not an assumption]
\label{thm:consistency-floor}
Let $(u^\star,p^\star,\nu_t)$ solve the steady incompressible RANS system on the open fluid
region $\Omega_f\subset\Omega$ of \autoref{sec:rf_operator}, with a \emph{linear
(Boussinesq) eddy-viscosity closure}, in which the
modelled deviatoric Reynolds stress is $-2\nu_tS$ with $S=\tfrac12(\nabla u+\nabla u^\top)$,
and with the isotropic part absorbed into the modified pressure:
\begin{equation}
u_j\partial_ju_i+\partial_iP-\partial_j\!\left[\nu_{\mathrm{eff}}(\partial_ju_i+\partial_iu_j)\right]=0,
\qquad \partial_ju_j=0 .
\label{eq:rans}
\end{equation}
Let $K\subset\subset\Omega_f$ be compact with $u^\star\in C^4(K)$ and $\nu_t\in C^1(K)$.
Then:
\begin{enumerate}
\item[\textup{(a)}] \emph{(The omitted term, in closed form.)} At every grid node in $K$,
for the \emph{dimensional} residuals of \eqref{eq:monitored} evaluated on
$u^\star|_{\Omega_h}$ (the per-equation scalings are reinstated in (c)),
\begin{equation}
\begin{gathered}
r_c=O(h^2),\qquad r_i \;=\; T_i+O(h^2),\\[2pt]
T_i:=(\partial_j\nu_t)(\partial_ju_i+\partial_iu_j)=2(S\nabla\nu_t)_i ,
\end{gathered}
\label{eq:omitted-term}
\end{equation}
and in two dimensions, where incompressibility makes $S$ symmetric and \emph{traceless},
\begin{equation}
\begin{gathered}
|T| \;=\; 2\lambda_S\,|\nabla\nu_t| \;=\; \sqrt{2}\,\|S\|_F\,|\nabla\nu_t|
\quad\text{\emph{exactly}},\\[2pt]
\text{with }\lambda_S=\sqrt{(\partial_xu)^2+\tfrac14(\partial_yu+\partial_xv)^2}
\end{gathered}
\label{eq:2d-identity}
\end{equation}
the strain-rate eigenvalue (distinct from the no-slip weight $\lambda$ swept later in this
section).
\item[\textup{(b)}] \emph{(Non-degeneracy.)} $T(x)=0$ if and only if $S(x)=0$ or
$\nabla\nu_t(x)=0$. In particular no orthogonality hypothesis is required.
\item[\textup{(c)}] \emph{(The floor survives refinement.)} If
$\bigl|\{x\in K: S(x)\neq0 \text{ and } \nabla\nu_t(x)\neq0\}\bigr|>0$, then
\begin{equation}
\liminf_{h\to0}\;\bigl\|R_h(u^\star|_{\Omega_h})\bigr\|
\;\ge\;\Bigl(\tfrac{|K|}{|\Omega|}\Bigr)^{1/2}\frac{\ell}{U^2}
\left(\frac{1}{|K|}\int_K |T|^2\,dx\right)^{1/2}\;>\;0 .
\label{eq:floor-limit}
\end{equation}
where $|\Omega|$ is the area of the \emph{whole crop}, since $\|\cdot\|$ averages over the
whole grid; the factor $(|K|/|\Omega|)^{1/2}\le1$ is the price of scoring a norm over
$\Omega$ while bounding it only on $K$.
Hence $R_h$ is \textbf{inconsistent} with \eqref{eq:rans}: hypothesis \textup{(H2)} of
\autoref{thm:residual-floor} is a consequence, not an assumption, and no grid refinement
removes the floor.
\end{enumerate}
\end{theorem}

\begin{proof}
(a) Expand the stress divergence in \eqref{eq:rans} using
$\partial_j[\nu_{\mathrm{eff}}(\partial_ju_i+\partial_iu_j)]
=\nu_{\mathrm{eff}}\nabla^2u_i+\nu_{\mathrm{eff}}\partial_i(\partial_ju_j)
+(\partial_j\nu_{\mathrm{eff}})(\partial_ju_i+\partial_iu_j)$.
The middle term vanishes identically by incompressibility of the \emph{exact} solution,
which is where this theorem lives; the discrete divergence of the rasterised field is a
different quantity and is discussed separately below. Also
$\nabla\nu_{\mathrm{eff}}=\nabla\nu_t$ because the laminar viscosity is constant. Substituting into \eqref{eq:rans} gives
$u_j\partial_ju_i+\partial_iP-\nu_{\mathrm{eff}}\nabla^2u_i=T_i$, whose left-hand side is
exactly the continuum form of the monitored $r_i$ in \eqref{eq:monitored}. Central first
differences are $O(h^2)$ on $C^3$ fields and the compact second difference is $O(h^2)$ on
$C^4$ fields, uniformly on the compact $K$; the products with the bounded coefficients
$u,v,\nu_{\mathrm{eff}}$ preserve that order. Continuity is a consistent discretisation of
$\nabla\!\cdot\!u^\star=0$, so $r_c=O(h^2)$ with no $O(1)$ term.
For \eqref{eq:2d-identity}: in two dimensions incompressibility gives $\partial_yv=-\partial_xu$,
so $S=\left(\begin{smallmatrix}a&b\\ b&-a\end{smallmatrix}\right)$ with $a=\partial_xu$,
$b=\tfrac12(\partial_yu+\partial_xv)$. Then $\det S=-(a^2+b^2)$ and the eigenvalues are
$\pm\lambda_S$, $\lambda_S=\sqrt{a^2+b^2}$. As $S$ is symmetric its singular values are
$|{\pm}\lambda_S|=\lambda_S$, both equal, so $|Sx|=\lambda_S|x|$ for \emph{every} $x$; apply
this to $x=\nabla\nu_t$ and use $\|S\|_F=\sqrt{2}\lambda_S$.
(b) Immediate from $|T|=2\lambda_S|\nabla\nu_t|$: a product of non-negative scalars vanishes
iff a factor does, and $\lambda_S=0$ iff $S=0$.
(c) On $K$ the non-dimensionalised integrand converges uniformly to
$(\ell/U^2)^2|T|^2$ by (a), so the Riemann sums of the mean square over $K\cap\Omega_h$
converge to the stated integral. Writing $n_\Omega$ for the number of grid cells and $n_K$
for those in $K$, and using that every discarded cell contributes non-negatively,
\[
\mathrm{MS}_\Omega=\frac{1}{n_\Omega}\sum_{\Omega_h}|R_h|^2
\;\ge\;\frac{1}{n_\Omega}\sum_{K\cap\Omega_h}|R_h|^2
\;=\;\frac{n_K}{n_\Omega}\,\mathrm{MS}_K ,
\]
and $n_K/n_\Omega\to|K|/|\Omega|$ as $h\to0$ since the grid is uniform on the crop. Taking
square roots and then $\liminf$ gives \eqref{eq:floor-limit}. Positivity follows from (b)
and the positive-measure hypothesis.
\end{proof}

\paragraph{Where the hypothesis fails, and why that is reassuring.}
\autoref{thm:consistency-floor}(b) says the floor degenerates exactly when $S=0$ or
$\nabla\nu_t=0$. Both hold for the uniform freestream, which is precisely the field that
\autoref{thm:residual-floor}(i) shows has an exactly zero monitored residual. The theorem
therefore predicts its own exception, and the exception is the spurious minimum. The
non-degeneracy hypothesis is also weak and checkable rather than cosmetic: it asks only
that the flow deform and the eddy viscosity vary somewhere in $K$, which is the definition
of a turbulent boundary layer or wake.

\paragraph{What this proves, and---emphatically---what it does not.}
It proves that the floor is bounded away from zero as $h\to0$: the ``refine the grid''
objection is closed \emph{in principle}. It does \textbf{not} explain the floor's measured
size, and we do not claim it does. Measured on the same fields with the same stencils, the
omitted term $T$ is $0.0016\to0.0048$ across the refinement ladder against a floor of
$0.0459\to0.1074$---about $3.5$--$4.8\%$---and \emph{repairing the full symmetric stress
divergence moves the floor by under $0.1\%$} ($0.10737\to0.10740$ at $512^2$;
these four figures come from the $16$-case, five-rung study scored on the
$0.1$-chord band, so they are not interchangeable with the $24$-case figures of
\autoref{sec:floor_ladder}, which use a different exclusion radius; the two agree on
the trend and on this fraction, which is what the argument uses;
\texttt{results/certificates/floor\_resolution\_decomposition.json}). What sets the size is
the mismatch between our Cartesian stencil and the body-fitted discrete balance the
reference solver actually closed, and the evidence for that is
\autoref{sec:floor_ladder}: continuity and momentum rise at the same rate in the same
cases, which no single-term story predicts---a closure omission moves momentum alone, and
lost divergence-freedom moves continuity alone. The two results are complementary and
neither substitutes for the other: \textbf{the theorem closes refinement in principle, the
ladder closes it in practice}, and the theorem's term is a small, separately provable
component that guarantees the floor cannot be refined away even if every other source were
eliminated. The algebraic identity of \eqref{eq:omitted-term} is verified numerically as a
committed unit test, converging at order $1.93$.

\paragraph{Why our own method-of-manufactured-solutions control cannot see this.}
The refinement study gates on a manufactured potential flow and measures order
$p=2.06$ in the interior band, which confirms the stencils are second-order and might be
read as confirming the operator is consistent. It cannot: for a potential flow
$u=\nabla\phi$ one has $\nabla^2u\equiv0$, so the monitored viscous term vanishes
identically and the manufactured field is an exact solution of the monitored operator
\emph{for any} $\nu_t$. The same artifact records this directly---multiplying $\nu_t$ by
$50$ changes the monitored norm by $6.7\times10^{-6}$ relative. An MMS gate built on a
field that annihilates the very term at issue measures truncation and is structurally
blind to inconsistency. We report the gate because it certifies the discretisation, and we
state its blind spot because a reviewer would otherwise supply it.

\subsection{The truth is not a minimiser, and residual descent leaves it}

\begin{theorem}[Residual floor: misplaced minimum, displaced minimiser, local divergence]
\label{thm:residual-floor}
Assume \textup{(H1)} $R_h$ is the monitored discrete operator of \eqref{eq:monitored}
(continuity $+$ momentum, no-slip omitted); and \textup{(H2)} the floor is nonzero,
$r^\star:=R_h(u^\star)\neq0$---which \autoref{thm:consistency-floor} supplies, and
\autoref{sec:floor_ladder} measures. Write $L:=DR_h(u^\star)$, let $\sigma_{\max}$ be its
largest singular value and $P:=P_{\operatorname{range}L}$. No injectivity is assumed; $L$
has a large kernel (\autoref{prop:kernel}). Then:
\begin{enumerate}
\item[\textup{(i)}] \emph{(Misplaced minimum, exact, and at every $h$.)} The
uniform-freestream field $u_\infty$ satisfies $R_h(u_\infty)=0$ exactly, so $\min_uJ=0$,
whereas $J(u^\star)=\tfrac12\|r^\star\|^2>0$. Hence $u^\star$ is \textbf{not} a minimiser
of $J$. This cannot be an artifact of the rasterisation or of the grid: a spatially
constant field has identically zero finite differences at \emph{every} $h$ and for
\emph{every} stencil in \eqref{eq:monitored}, one-sided edges included.
\item[\textup{(ii)}] \emph{(Misplaced stationarity.)} $\nabla J(u^\star)=L^\top r^\star$,
so $u^\star$ is stationary for $J$ iff $L^\top r^\star=0$; if $r^\star\notin\ker L^\top$,
gradient flow on $J$ leaves $u^\star$.
\item[\textup{(iii)}] \emph{(Displacement of the minimiser, bounded below.)} The
minimisers of the linearised objective $\tfrac12\|r^\star+Le\|^2$ form the affine set
$e_\infty+\ker L$ with $e_\infty=-L^{+}r^\star$ ($L^{+}$ the Moore--Penrose
pseudoinverse). \textbf{Every} point $e$ of that set satisfies
\begin{equation}
\|e\|\;\ge\;\|e_\infty\|\;\ge\;\frac{\|P\,r^\star\|}{\sigma_{\max}},
\label{eq:rf-lower}
\end{equation}
which is strictly positive whenever $r^\star$ has any component in $\operatorname{range}L$.
\item[\textup{(iv)}] \emph{(Local divergence direction.)} Under residual-gradient flow
$\dot e=-\nabla J(\hat u)$,
$\frac{d}{dt}\tfrac12\|e\|^2\big|_0=-\|Le\|^2-(Le)\cdot r^\star$, which is positive on the
open cone $(Le)\cdot r^\star<-\|Le\|^2$. The cone is non-empty whenever $Pr^\star\neq0$,
and is reached by $Le=-\alpha\,Pr^\star/\|Pr^\star\|$ for $0<\alpha<\|Pr^\star\|$.
\end{enumerate}
\end{theorem}

\begin{proof}
(i) Every stencil in \eqref{eq:monitored} annihilates a constant---central differences
term-by-term, and the one-sided edge stencils $(f_1-f_0)/h$ and $(f_0-2f_1+f_2)/h^2$ by
inspection---so $R_h(u_\infty)=0$ identically, giving $J(u_\infty)=0=\min J$ since $J\ge0$.
By (H2), $J(u^\star)>0$.
(ii) $\nabla J(u)=DR_h(u)^\top R_h(u)$; evaluate at $u^\star$ and use \eqref{eq:rf-decomp}.
(iii) $\tfrac12\|r^\star+Le\|^2$ is convex, and its minimiser set is the solution set of the
normal equations $L^\top Le=-L^\top r^\star$, namely $e_\infty+\ker L$ with
$e_\infty=-L^{+}r^\star$; no injectivity is needed. Since $L^{+}$ maps into
$\operatorname{range}(L^\top)=(\ker L)^\perp$, any minimiser is $e=e_\infty+k$ with
$k\in\ker L$ orthogonal to $e_\infty$, so $\|e\|^2=\|e_\infty\|^2+\|k\|^2\ge\|e_\infty\|^2$:
the bound holds uniformly over the minimiser set, not merely at the minimum-norm
tie-break. For the second inequality, $LL^{+}=P$ gives $Le_\infty=-Pr^\star$, whence
$\|Pr^\star\|=\|Le_\infty\|\le\sigma_{\max}\|e_\infty\|$.
(iv) To first order $\dot e|_0=-L^\top(r^\star+Le)$, so
$\frac{d}{dt}\tfrac12\|e\|^2|_0=-e^\top L^\top(r^\star+Le)=-\|Le\|^2-(Le)\cdot r^\star$,
positive iff $(Le)\cdot r^\star<-\|Le\|^2$. Since $Le$ is confined to
$\operatorname{range}L$, only the \emph{projected} floor is reachable, which gives the
constant $\|Pr^\star\|$ and not $\|r^\star\|$.
\end{proof}

\paragraph{Scope, stated before a reviewer states it for us.}
Legs (iii) and (iv) are statements about the \emph{linearised} objective at $u^\star$ and
about gradient flow; the deployed corrector is supervised toward truth and the acceptance
gate uses the residual as a one-bit veto, so neither leg describes the deployed loop. They
describe residual descent---which this paper does run, from the exact truth and with no
network anywhere in the loop, and which behaves as (ii) and (iv) predict:
$\nabla J(u^\star)\neq0$ while $\nabla\|e\|^2(u^\star)=0$ exactly, and $300$ steps cut the
residual by $84\%$ in $24/24$ cases while taking the field error from zero to a median of
$0.91$ (\autoref{sec:descent}). The nonlinear, non-asymptotic version of (iii)--(iv) is
that measurement; the theorem is the local mechanism. We also do not report a number for
\eqref{eq:rf-lower}: $\sigma_{\max}$ and $\|Pr^\star\|$ are not measured here, so (iii) is
a strictly-positive-displacement result and we do not describe it as a quantified
detection limit.

\subsection{The kernel: what the monitor cannot see at all}

The previous version of this section offered two examples of null modes. The situation is
considerably worse than two examples, and it is provable by counting.

\begin{proposition}[Structure of the undetectable subspace]
\label{prop:kernel}
With $M$, $W$, $A=M\setminus W$ and $L:\mathbb{R}^{4|M|}\to\mathbb{R}^{3|A|}$ as above:
\begin{enumerate}
\item[\textup{(K1)}] \emph{(Counting.)}
$\dim\ker L\;\ge\;4|M|-3|A|\;=\;|M|+3|W|\;\ge\;|M|$. At least \textbf{one quarter} of the
fluid state space is invisible to the monitor at first order, for the structural reason
that $R_h$ imposes three equations per cell on four unknowns per cell, plus three further
dimensions for each wall-ring cell whose residuals are zeroed outright.
\item[\textup{(K2)}] \emph{(Solid degrees of freedom, exactly and nonlinearly.)} If the
body is at least three cells from the outer border of the crop, then every degree of
freedom on a solid cell lies in $\ker R_h$ \emph{exactly}---not merely in $\ker L$:
$R_h$ is unchanged by \emph{any} perturbation supported on the solid, of any size.
\item[\textup{(K3)}] \emph{(The $\nu_t$ block is diagonal, in closed form, exactly.)}
$R_h$ is \emph{affine} in $\nu_t$, which enters only as the pointwise multiplier
$-\nu_{\mathrm{eff}}\Delta_hu$. Hence for any perturbation $\delta$ supported on the
$\nu_t$ channel with $\nu+\nu_t+\delta\ge0$,
\begin{equation}
R_h(u^\star+\delta)-R_h(u^\star)
=\bigl(0,\;-\delta\odot\Delta_hu^\star,\;-\delta\odot\Delta_hv^\star\bigr)\big|_A
\quad\text{exactly},
\end{equation}
so $\|R_h(u^\star+\delta)-R_h(u^\star)\|^2=\sum_{k\in A}c_k^2\delta_k^2$ with
$c_k=\frac{\ell}{U^2}\sqrt{(\Delta_hu^\star_k)^2+(\Delta_hv^\star_k)^2}$. Consequently, for
$\Omega_\varepsilon:=\{k\in A:c_k\le\varepsilon\}$ and any $\delta$ supported in
$\Omega_\varepsilon$, the monitor moves by at most $\varepsilon\|\delta\|$: the sublevel
set $\{u:\|R_h(u)\|\le t\}$ contains a ball of radius $(t-\|r^\star\|)/\varepsilon$ inside
a coordinate subspace of dimension $|\Omega_\varepsilon|$, for every $t>\|r^\star\|$.
\item[\textup{(K4)}] \emph{(A Nyquist pressure mode the monitor barely registers.)} Let
$p_\varepsilon(i,j)=\varepsilon(-1)^{i+j}$. Then $D_xp_\varepsilon=D_yp_\varepsilon=0$ at
every \emph{interior} node, so an $O(1)$ checkerboard oscillation in the pressure changes
the monitored residual at \emph{no} active cell except those on the outer border ring of
the crop---an $O(1/N)$ fraction of the active set---and would be \emph{exactly} invisible
were that ring excluded, as the wall ring already is.
\end{enumerate}
\end{proposition}

\begin{proof}
(K1) Rank--nullity, with $\operatorname{rank}L\le3|A|$ and $|A|=|M|-|W|$; the final
inequality is $|M|+3|W|\ge|M|=\tfrac14\cdot4|M|$.
(K2) Every interior stencil in \eqref{eq:monitored} reads only a cell's four neighbours
($D_x,D_y$ reach $(i\pm1,j)$; the compact $\Delta_h$ reaches $(i\pm1,j)$ and $(i,j\pm1)$).
By definition of $W$, no active cell has a solid $4$-neighbour, so no active cell's stencil
reads a solid value. The one-sided edge stencils reach two cells inward, which is why the
three-cell separation from the outer border is required; it holds for the $3c\times3c$ crop
about a unit chord used here. The statement is nonlinear because $R_h$ simply never
evaluates those entries.
(K3) $\nu_{\mathrm{eff}}=\nu+\nu_t$ multiplies $\Delta_hu$ pointwise and appears nowhere
else in \eqref{eq:monitored}, so $R_h$ is affine in $\nu_t$ with the stated diagonal linear
part; the clip $\max(\nu_{\mathrm{eff}},0)$ is inactive under the stated hypothesis, which
holds on the reference data since $\nu_t\ge0$ there.
(K4) $(-1)^{(i+1)+j}=(-1)^{(i-1)+j}$, so the central difference
$(p_{i+1,j}-p_{i-1,j})/2h$ vanishes identically, and likewise in $y$; only the first-order
one-sided stencils at the four borders return a nonzero value, $2\varepsilon/h$.
\end{proof}

\paragraph{Reading the kernel.} (K3) has a physical statement: \emph{eddy viscosity is
unobservable from a momentum residual wherever there is nothing to diffuse}. In the
freestream $\Delta_hu^\star\approx0$, so $c_k\approx0$ and $\nu_t$ may be set arbitrarily
there without moving the monitor at all. This is why we never claim the residual certifies
the $\nu_t$ channel. (K4) is the classical odd--even (checkerboard) decoupling of pressure
on a collocated grid with central differencing, whose standard remedy is momentum
interpolation \citep{rhiechow1983}; a surrogate-side monitor assembled from plain central
differences inherits the pathology without the remedy, and we flag it rather than claim the
checkerboard as an exact null mode---the border stencils mean it is not one.

\subsection{Ranking survives; certification does not}

\paragraph{Why the floor defeats certification but not triage.} Certification asks for a
map from $\|R_h(\hat u)\|$ to a bound on $\|\hat u-u^\star\|$. It fails here for a reason
that needs no linearisation: $u^\star$ itself is rejected at every threshold below
$\|r^\star\|$ (\autoref{thm:residual-floor}(i)), the sublevel sets are unbounded in the
directions of \autoref{prop:kernel}, and the floor sits \emph{above} the operating point
(mean $\|r^\star\|=0.192$ versus $\|R_h(\hat u)\|=0.114$). Ranking asks something much
weaker and \emph{between} cases: that a case whose prediction is worse tend to carry a
larger monitored residual. Nothing above forbids that, and it is measured---$\rho=0.61$
and AUROC $0.871$ on field error, $0.952$ on worst-decile drag error
(\autoref{sec:selective}). We therefore make the ranking claim empirically and do
\emph{not} derive it from \eqref{eq:rf-decomp}. In particular we do not argue it from a
two-sided bound $\sigma_{\min}\|e\|\le\|Le\|\le\sigma_{\max}\|e\|$: $\sigma_{\min}$ is the
smallest nonzero singular value of a $\sim\!4.7\times10^4$ by $6.6\times10^4$
advection--diffusion Jacobian with a continuum of near-null directions, it is not
separable from zero numerically, and a bound resting on it would be vacuous. Equation
\eqref{eq:rf-decomp} offers only the heuristic that far from the truth $\|Le\|$ dominates
$\|r^\star\|$; that is a regime intuition, not a guarantee, and the guarantee we do ship
for the error bound is the distribution-free conformal one, not a residual bound.

\paragraph{The apparent contradiction, and its resolution.} If the deployed prediction sits
\emph{below} the floor in $80\%$ of cases, it sits in the regime the decomposition calls
floor-dominated---yet $\rho=0.61$ is measured there. The two are statements about
different things. The decoupling is \emph{within} a case: holding the case fixed, moving
the field to lower residual need not move it toward the truth, which is exactly what
\autoref{sec:descent} measures. The detector claim is \emph{between} cases, where both
$r^\star$ and $Le$ vary. That raises the obvious worry---that the ranking is a difficulty
proxy rather than an error detector---and we test it: conditioning on each case's own floor
leaves a partial rank correlation of $0.561$ against a raw $0.610$
(\autoref{sec:selective}). The detector survives the control, and should be read as
between-case triage, never as evidence that residual descent within a case moves toward the
truth.

\paragraph{Empirical confirmation.} On $200/200$ real AirfRANS test cases the monitored
residual of the ground-truth field is substantially nonzero: $\|r^\star\|$ has mean $0.192$
(median $0.133$) when averaged over fluid cells, the convention of the cited artifact, and
$0.191$ in the whole-grid norm named above---the two differ only by the geometric factor
$(|\Omega_f|/|\Omega|)^{1/2}$, here $0.995$. Meanwhile the uniform-freestream field
gives $\|R_h(u_\infty)\|=0$ in every case---the objective strictly prefers a physically
wrong field to the truth (\texttt{results/certificates/residual\_floor\_realdata.json}).
Because every field on a given case is scored against the same mask, that factor cancels in
any comparison, so the counts below are identical under either convention. A
trained backbone's prediction (the dropout-FNO of
\texttt{checkpoints/certificates\_deq.pt}, distinct from the Transolver headline) sits
\emph{below} the truth floor in $160/200$ cases. We give that number with its
counterweight rather than on its own, because it is the flattering one: on the
\emph{deployed} Transolver the same inversion occurs on only $2.5$--$7.0\%$ of cases
across seeds ($8/5/5$ of $200$ for the backbone alone and $12/9/14$ for
Transolver$+$DEQ, seeds $0,1,2$;
\texttt{results/certificates/transolver\_inversion.json},
\texttt{results/review/functional\_audit\_gate\_analysis.json}), so the
truth is the best-scoring field on more than $93\%$ of them, and the median gap reverses
sign ($+0.013$ to $+0.015$ \emph{above} the floor, against $-0.024$ below it for the
dropout-FNO). Both backbones are scored with an identical ruler: the harness reproduces
the published dropout-FNO row to all digits. The inversion is a property of that one
backbone and we rest no claim on it. We also decline the obvious explanation, because it
does not survive measurement: the inversion is \emph{not} simply over-smoothing. The
accurate backbone is not smoother than the truth---its velocity-gradient energy is
$1.00\times$ the truth's in a near-body band---so smoothing cannot be what separates the
two. (The dropout-FNO measures \emph{rougher} still, $1.75\times$ over the fluid, though we
lean on that number only lightly: it is a grid-native field whereas the Transolver field is
rasterised from the native cloud, so the cross-backbone gradient comparison carries a
rasterisation confound that the inversion counts themselves do not.) What does separate
them is that the dropout-FNO's monitored residual is nearly case-\emph{independent} (spread
$0.034$ versus the truth's $0.162$ and the Transolver's $0.161$), so it neither tracks case
difficulty nor clears the floor. That is what \eqref{eq:rf-decomp} predicts: a prediction
falls below the floor only when its error is systematically anti-aligned with $r^\star$,
whereas unbiased error of any magnitude \emph{raises} the monitored residual---which is why
the accurate backbone sits above it. What does not depend
on the backbone is the floor itself, which no model enters. The floor is not an artifact of the
deployed grid: re-rasterised from the source point clouds at $128^2$--$512^2$ and scored on
a region fixed in physical units it \emph{increases} by $2.6\times$, fitting
$\|r^\star\|\sim h^p$ with $p=-0.64\pm0.29$ and decaying in $0/16$ cases, while a
manufactured potential flow pushed through the identical rasterisation converges at order
$2.02\pm0.05$ and a $C^1$ re-rasterisation \emph{raises} rather than lowers the floor
(\autoref{sec:floor_ladder}). Those band-restricted magnitudes are a different scoring
region from the $0.192$ above and are not comparable to it; the refinement claim is a
statement about the trend within one scoring policy.

\paragraph{Robustness to the boundary term.} One might object that the monitored $R_h$
omits the no-slip term, so the negative is merely an artifact of dropping the boundary
constraint. We tested this with the \emph{framework's own} boundary-inclusive residual
$\rho_{\rm res}=\sqrt{r_c^2+r_x^2+r_y^2+r_{bc}^2}$ (the trust-map definition, $r_{bc}$ the
proximity-weighted no-slip penalty) along the correction sweep (dropout-FNO$+$DEQ, $n{=}10$,
same regime as \autoref{tab:iters}; \texttt{results/control/bc\_inclusive\_sweep.json}). The
negative \emph{survives}: as the corrector cuts field error ($\texttt{mse\_u}$
$3.17\!\to\!2.32$), the boundary-\emph{inclusive} residual rises ($0.145\!\to\!0.811$) in
lockstep with the boundary-excluded one ($0.125\!\to\!0.807$)---along the correction path the
prediction and the truth both approximately satisfy no-slip, so $r_{bc}$ stays small and
roughly constant and the interior floor dominates. The objection generalises---\emph{some}
no-slip weight might rescue the residual---so we sweep it. Writing the weighted monitor
$\rho_\lambda=\sqrt{\|R_h\|^2+\lambda\,\overline{r_{bc}^2}}$, both logged components are
recoverable at every sweep point, so $\lambda$ can be swept post hoc at no computational
cost (\texttt{scripts/bc\_weight\_sweep.py}, \texttt{results/control/bc\_weight\_sweep.json}).

\textbf{(A)}~\emph{At the deployed resolution} the uniform field remains a spurious
minimum: it has both a smaller interior residual and a smaller no-slip term than the truth
($\overline{r_{bc}^2}=0.0053$ versus $0.0073$; the ordering holds in $14/16$ cases of the
resolution study), hence $\rho_\lambda(u_\infty)<\rho_\lambda(u^\star)$ for all
$\lambda\ge0$ \emph{at that resolution}. We claim no more, and state the limitation rather
than the closed form we previously claimed for every grid: this leg is
resolution-contingent. The framework's proximity weight decays over $3\min(dx,dy)$ and so
scores a physically \emph{thinner} band as the grid refines, and the margin collapses
monotonically---$+9.3\%$ at $128^2$, $+7.4\%$ at $181^2$, $+6.9\%$ at $256^2$, $+3.9\%$ at
$362^2$, and $-2.8\%$ at $512^2$, where it has crossed. Where it crosses, the crossover
weight
$\lambda^\star=\|R_h(u^\star)\|^2/(\overline{r_{bc}^2}(u_\infty)-\overline{r_{bc}^2}(u^\star))$
has median ${\approx}600$ (range $359$--$8549$). Since the deployed trust map uses
$\lambda=1$, a monitor that rescued the ordering would be one in which the no-slip term
outweighs the entire PDE residual by two orders of magnitude---a real technical break, but
not a practical rescue, because such an object is no longer a PDE monitor. Two further
facts close off readings of (A) that would otherwise be live: the far-field ring
contributes under $1\%$ of $\overline{r_{bc}^2}(u^\star)$, so the ordering is genuinely
about the near-wall band and not an artifact of the crop's outer boundary; and the
ordering was never universal, holding at all five rungs in only $11/16$ cases.

\textbf{(B)}~Along the correction path, where field error falls ($3.17\!\to\!2.32$),
\emph{both} increments are positive---the interior term by $+0.635$ and the no-slip term by
$+0.0019$---so $\rho_\lambda$ rises for every $\lambda\ge0$; there is no crossover weight.
Leg (B) does not depend on the truth-versus-uniform ordering of (A), and it is leg (B) that
carries the argument for \autoref{tab:iters}: up-weighting the boundary condition cannot
convert this residual into a usable correction objective, because the no-slip violation
grows along the very path that reduces the error.

\paragraph{Assumptions, stated plainly.} \autoref{thm:residual-floor} is about the
\emph{monitored discrete} operator $R_h$, not the continuous PDE.
\autoref{thm:consistency-floor} is the bridge between them, and it is where the smoothness
hypotheses live: it holds on a compact subset of the open fluid domain on which the
solution is $C^4$ and $\nu_t$ is $C^1$, which excludes a neighbourhood of the wall---the
same methodological restriction the refinement ladder makes when it scores a region fixed
in physical units. (H2) therefore holds here for a reason that is \emph{not} grid
resolution, and we have retired the earlier claim that it followed from the $128^2$ grid
under-resolving the boundary layer. It follows from operator provenance: the floor grows
under refinement, the individually grid-converged convective and pressure-gradient terms
($0.183\!\to\!0.199$ and $0.180\!\to\!0.181$) fail to cancel under a discretisation
different from the one that produced the labels, and restoring the omitted stress
divergence moves the floor by under $0.1\%$. One further honest bound: the measured
$r^\star$ is evaluated on a \emph{rasterised} $u^\star$---a linear interpolant of an
unstructured finite-volume solution---which is not $C^4$ and carries its own
representation error, so the measured floor mixes the theorem's term with that error. That
is why the theorem is stated as a limit on the exact solution and the magnitude is reported
as a measurement, and why we never present one as evidence for the other.

\paragraph{Relation to prior work.} \citet{krishnapriyan2021failure} tie ill-conditioning
to the \emph{optimisation difficulty} of PINN losses, and \citet{wang2022ntk} analyse
spectral bias and loss-term imbalance in PINN \emph{training dynamics}; both concern
trainability, not a displacement of the residual minimiser from the truth. The
forward-versus-backward error gap is classical \citep{trefethenbau1997}, and a-posteriori
PINN bounds upper-bound error \emph{by} residual---the benign regime, and one calibrated by
the identity $R(u^\star)=0$, which is exactly the premise our monitor lacks.

The sharpest way to state where the novelty sits is through goal-oriented estimation.
Dual-weighted-residual methods \citep{beckerrannacher2001dwr} exist precisely because
residual magnitude is not an error estimate, and they work beautifully for reduced-order
models---because there the residual is $r=A_{\rm FOM}(u_{\rm ROM})-b$, which vanishes at
the truth \emph{by construction}, since the full-order operator is available and is the one
whose root the reference is. A dataset-trained surrogate has no $A_{\rm FOM}$. The
obstruction we report is therefore about \textbf{operator availability}, not about RANS,
not about two dimensions and not about the grid: it appears whenever the operator one can
afford to evaluate at deployment is not the operator that generated the labels. We do not
claim that a dual-weighted estimate provably fails here---that would require constructing
the adjoint $z$ and bounding $\langle z,r^\star\rangle$ from below, which we have not
done---only that the identity it relies on acquires an uncontrolled $\langle z,r^\star\rangle$
term, and that forming the adjoint needs a linearised solve on the deployed grid, which is
the cost the surrogate exists to avoid.

The nearest classical relatives are worth naming precisely, because each has the structure
we need and none has the case. Defect correction \citep{stetter1978defect} and multigrid
$\tau$-correction \citep{brandt2011multigrid} both drive one operator using the residual of
\emph{another}, and both handle the resulting mismatch by carrying an explicit correction
term---exactly the object a deployment-time monitor cannot form, since it would require the
fine operator we are trying to avoid running. \citet{zhang2026phymgn} reach the same
practical conclusion from the data side for shock physics, and their remedy---defining the
loss relative to the residual present in the ground-truth data---likewise requires
$R_h(u^\star)$, which a deployment-time monitor by definition does not have. In the
finite-element a-posteriori literature, an error contribution the estimator cannot see and
that does not vanish under refinement is \emph{data oscillation} \citep{morin2000data}; our
floor is its analogue for an inconsistent monitored operator, with the difference that data
oscillation is a property of the right-hand side's representation while ours is a property
of the operator itself. Least-squares finite element methods
\citep{bochevgunzburger2009book} obtain a residual functional that \emph{is}
norm-equivalent to the error---the property our monitor lacks---by constructing the
discretisation and the functional together, which is precisely the option a surrogate-side
audit forgoes when it adopts a cheap operator after the fact. And where the monitored
operator \emph{is} the solver's own, residual-driven correction works on steady CFD
\citep{lei2026newtonkrylov}, which locates our negative rather than contradicting it.

That a residual blind to certain modes makes the truth a non-minimiser is, in its bare
form, elementary---the pressure gauge above. Our contribution is threefold and narrower
than earlier drafts of this work claimed. \textbf{(a)} A \emph{proof} that this floor has a
nonzero continuum limit (\autoref{thm:consistency-floor}), so it is a property of operator
provenance rather than of the mesh and cannot be refined away even in principle---the
$\nu_{\mathrm{eff}}\nabla^2u$ simplification of the RANS momentum equation is the standard
choice in physics-informed CFD surrogates, so this is a floor an entire practice has not
accounted for. We previously advertised instead an upper bound
$\|e_\infty\|\le\|r^\star\|/\sigma_{\min}$ as ``the quantified floor''; that inequality
points in the wrong direction for the thesis and is numerically vacuous at this problem
size, and we withdraw it. \textbf{(b)} An exact characterisation of the undetectable
subspace (\autoref{prop:kernel})---a counting bound of at least a quarter of the fluid
state, exact nonlinear invisibility of the solid, a diagonal $\nu_t$ block with closed-form
entries, and the Nyquist pressure mode---replacing what was previously an existence claim
with two examples. \textbf{(c)} The demonstration on a learned neural-operator CFD
corrector on AirfRANS that the resulting monitor can rank predictions, cannot be
descended, and certifies them only at a width the floor rather than the model sets.
